
Gradient-based interventions for spurious correlation robustness offer the possibility of single-run training without group labels. However, experiments conducted in this work reveal a measurement limitation: gradient subspace accumulation with one sample per epoch produces representations too low-rank to distinguish spurious from core feature directions. Where theory predicted clear separation (greater than 70\% spurious alignment versus less than 30\% core alignment based on simplicity bias), both alignments measured approximately 5\%---statistically indistinguishable.

This finding is relevant because methods claiming to exploit training dynamics for robustness may measure noise rather than meaningful signal without understanding minimum requirements for gradient subspace analysis.

\subsubsection{1.1 The Problem}

Deep neural networks trained with empirical risk minimization (ERM) learn spurious correlations---shortcuts that achieve low training loss but fail on minority groups. On Waterbirds, ERM models learn to associate waterbirds with water backgrounds (95\% correlation in training data), achieving approximately 60\% worst-group accuracy when waterbirds appear on land backgrounds.

Existing solutions require resources unavailable in many practical settings. Group DRO requires group annotations during training. Just Train Twice (JTT) requires two complete training runs. Deep Feature Reweighting (DFR) requires held-out data with group labels for last-layer retraining. The possibility of training dynamics exploitation---intervening on gradients during a single training run without labels---remains unvalidated.

\subsubsection{1.2 Investigation}

This work tested the foundational assumption underlying Progressive Gradient Orthogonalization: that early gradient subspaces primarily capture spurious feature directions. ResNet-50 was trained on Waterbirds, gradients were accumulated during epochs 1-10, a top-50 SVD subspace was computed, and alignment was measured with spurious (background-varying) and core (bird-type-varying) gradient directions.

The result was unexpected. Both alignments registered at approximately 0.05---far below the 0.70 and 0.30 thresholds specified by the experimental design, and statistically indistinguishable from each other. Investigation revealed the cause: accumulating one gradient per epoch (the last batch only) produced a 10-dimensional subspace in a 25.6-million-parameter space. This rank is insufficient to capture meaningful variance in any direction. The measurement apparatus failed before the hypothesis could be tested.

\subsubsection{1.3 Contributions}

This investigation yields a methodological constraint rather than a positive result:

\begin{enumerate}
\item Gradient subspace analysis in high-dimensional parameter spaces requires sufficient sample density---single gradients per epoch produce subspaces too low-rank for meaningful directional analysis.
\end{enumerate}

\begin{enumerate}
\item Minimum requirements for valid gradient subspace measurement are documented: multi-batch accumulation within early epochs is necessary, not single-batch per epoch sampling.
\end{enumerate}

\begin{enumerate}
\item A template is provided for future gradient-based robustification work to validate measurement apparatus before testing hypotheses.
\end{enumerate}

The Progressive Gradient Orthogonalization hypothesis remains unverified. What is established is what valid testing requires.

